\documentclass[DIN, pagenumber=false, fontsize=11pt, parskip=half]{scrartcl}
\usepackage[UTF8]{ctex}  % 中文支持
% \usepackage{ngerman}
\usepackage[utf8]{inputenc}
\usepackage[T1]{fontenc}
\usepackage{textcomp}
\usepackage{xcolor} % Required for custom red

\usepackage{geometry} % Custom margins

\usepackage{tocloft} % Modifying the table of contents

% for matlab code
% bw = blackwhite - optimized for print, otherwise source is colored
\usepackage[framed,numbered,bw]{mcode}
\usepackage[most]{tcolorbox}
% for other code
%\usepackage{listings}

\setlength{\parindent}{0em}

% set section in CM
\setkomafont{section}{\normalfont\bfseries\Large}

\newcommand{\mytitle}[1]{{\noindent\Large\textbf{#1}}}
\newcommand{\mysection}[1]{\textbf{\section*{#1}}}
\newcommand{\mysubsection}[2]{\romannumeral #1) #2}

\newtcolorbox{parchmentquote}{
    enhanced,
    colback=orange!5!white,   % 浅米黄色背景
    colframe=orange!50!black,
    boxrule=0.5pt,
    arc=2pt,                  % 细微圆角
    boxsep=5pt,
    before skip=10pt,
    after skip=10pt,
    fontupper=\itshape,
    borderline west={2pt}{0pt}{orange!80!black} % 加厚的左装饰线
}
%===================================
\begin{document}

\noindent\textbf{课下练习} \hfill \textbf{online course}\\
primary school \hfill Liu\\

\mytitle{简单代数式 \hfill \today}
\newline
\begin{parchmentquote}
	这一部分就是小学比较难的内容了，需要慢慢消化，不是一两节课就能掌握的，而是一种思维的转变。
\end{parchmentquote}


%===================================
\section{代数式}

\subsection{用字母表示数} 
把数换成字母，来表示一些数量关系。
如：

\begin{enumerate}
    \item 一支铅笔$a$元
    \item 小明今年$b$岁
\end{enumerate}
\subsection{书写规范} 
\begin{enumerate}
    \item 字母和数字或者字母与字母之间用$\cdot$连或者省略,如$a \cdot b$,$ab$,$5x$。
    \item 省略乘号时，数字在前，字母在后，如$5x$，而不是$x5$。
    \item 数字是$1$或者$-1$时，通常省略不写，如$a$,$-a$ 。
    \item 数字是带分数时候，要写成假分数形式，如$\frac{7}{3}a$，而不是$2\frac{1}{3}a$。 
    \item 运算结果不出现除号，写成分数形式。如$a \div b$,写成$\frac{a}{b}$，$a \div 2$,写成$\frac{a}{2}$。
\end{enumerate}


\section{练习}
 用字母表示下列数量关系：
\begin{enumerate}
  
    \item 已知小红的年龄是$a$岁，试用代数式表示下列关系：
    \begin{enumerate}
        \item 小明今年的年龄是小红的两倍。
        \item 小明今年的年龄比小红大3岁。
        \item 小明今年的年龄是小红的年龄的3倍再多2岁。
    \end{enumerate}
\end{enumerate}
尝试列方程，未知数数用$x$表示：
\begin{enumerate}
    \item 小明今年的年龄是小红的两倍，小明今年12岁，试列方程求小红的年龄。
    \item 小明今年的年龄比小红大3岁，小明今年12岁，试列方程求小红的年龄。
    \item 小明今年的年龄是小红的年龄的3倍再多2岁，小明今年12岁，试列方程求小红的年龄。  
    \item  小明今年的年龄是小红的两倍，小明今年12岁，试列方程求小红的年龄。
\end{enumerate}
\end{document}